\documentclass[10pt,a4paper]{amsart}
\usepackage[margin=19mm]{geometry}
\usepackage{amsmath,amssymb,mathtools}
\usepackage[scr=rsfs,cal=euler]{mathalfa}
\usepackage{xfrac}
\usepackage[unicode,hidelinks,bookmarks=false]{hyperref}
\newcommand{\ie}{i.e.\ }
\newcommand{\cf}{cf.\ }
\newcommand{\resp}{respectively\ }
\newcommand{\Subsection}{\subsection}
\providecommand{\nobreakdash}{\nobreak\hbox{-}}
\setlength{\emergencystretch}{2em}
\multlinegap0pt
\usepackage[shortlabels]{enumitem}
\usepackage{amssymb}
\usepackage{mathtools}
\usepackage{xcolor}
\newtheorem{lemma}{Lemma}
\title{On sporadic symmetry breaking operators for principal series representations of the de Sitter and Lorentz groups}
\author{V\'{i}ctor P\'{E}REZ-VALD\'{E}S}
\date{Source: arXiv:2506.23064v5 (CC BY 4.0)}
\begin{document}
\maketitle

\noindent\emph{Source and licence.} On sporadic symmetry breaking operators for principal series representations of the de Sitter and Lorentz groups. Version arXiv:2506.23064v5 (CC BY 4.0), distributed by arXiv under the Creative Commons Attribution 4.0 International licence, http://creativecommons.org/licenses/by/4.0/, as recorded in the arXiv licence metadata for this exact version and shown on the abstract page https://arxiv.org/abs/2506.23064v5. Full licence notice: \texttt{licenses/arxiv-2506.23064-CC-BY-4.0.txt}.\par\smallskip

\noindent\textit{Editorial scope.} Counted unit: the article's lemma lemma-P=Q (source0.tex lines 1545--1547). The article's complete original proof of it is delivered with it (source0.tex lines 1548--1612, one \texttt{proof} environment of 65 lines). No mathematics is written, repaired or re-derived: the statement and the proof are the article's own lines, in the article's own notation, and no retained line is substituted or re-flowed.  Retained uncounted as the source-local context the proof needs: lines 1233--1241; lines 1527--1532; lines 1534--1539; lines 1541--1543; lines 2328--2335; lines 2337--2342; lines 2394--2404.  Nothing was omitted from any retained span, so the package carries no omission record.  No retained line contains a citation, so the wrapper carries no bibliography and no bibliography entry.  No retained line contains a figure environment, an image-inclusion command or drawing code, so no image is delivered and the package declares no asset dependency.  Editorial formatting only, in addition: an AMS \texttt{amsart} preamble that restates the article's own declarations and macros used by the retained lines, namely the environments \texttt{lemma} on separate counters, together with the standard packages those lines need.  The retained environments print in this document as Lemma 1 in order of appearance, whereas the article prints the same items with the numbering of its own full document.  The complete native source of the article as distributed in the arXiv e-print for this exact version is bundled unchanged as \texttt{sources/180676/source0.tex}.
\begin{equation}\label{definition-constants-Cjr+-}
C_{j, r}^\pm := \begin{pmatrix}
N-j\\
r
\end{pmatrix}
\frac{\Gamma(2N-r+1)}{\Gamma(N+j+1)}
\frac{\Gamma(N\mp m+1)}{\Gamma(N\mp m -r +1)}
\frac{\Gamma(\lambda+a-N+r-1)}{\Gamma(\lambda+a-N-1)}.
\end{equation}

\begin{align}
\label{def-alpha}
\alpha(\lambda, N, a, m) &:= \sum_{r = 0}^{N}(-1)^r C_{0,r}^+\frac{\Gamma\left(\left[\frac{a+m}{2}\right]\right)}{\Gamma\left(\left[\frac{a+m}{2}\right]-N+r\right)},\\
\label{def-beta}
\beta(\lambda, N, a, m) &:= \sum_{r = 0}^{N}(-1)^r C_{0,r}^+\frac{\Gamma\left(\left[\frac{a+m+2}{2}\right]\right)}{\Gamma\left(\left[\frac{a+m+2}{2}\right]-N+r\right)},
\end{align}

\begin{equation}\label{def-P(lambda)}
\begin{aligned}
P(\lambda):= &\left(\lambda+\left[\frac{a+m-1}{2}\right]\right)\left[\frac{a+m}{2}\right]\alpha(\lambda, N, a, m)\beta(\lambda, N, a, -m)\\
&- \left(\lambda+\left[\frac{a-m-1}{2}\right]\right)\left[\frac{a-m}{2}\right]\alpha(\lambda, N, a, -m)\beta(\lambda, N, a, m).
\end{aligned}
\end{equation}

\begin{equation}\label{def-Q(lambda)}
Q(\lambda) := \frac{\Gamma(N+m+1)\Gamma(N-m+1)}{\Gamma(m)\Gamma(1-m)}\prod_{q=0}^{2N}(\lambda-N-1+a+q).
\end{equation}

\begin{equation}\label{def-hypergeometric}
{}_pF_q\left(
\begin{matrix}
a_1, \ldots, a_p\\
b_1, \ldots, b_q 
\end{matrix}
; z\right) := \sum_{n=0}^{\infty} \frac{(a_1)_n \cdots (a_p)_n}{(b_1)_n \cdots (b_q)_n}\frac{z^n}{n!},
\end{equation}

\begin{equation}\label{Pochhamer-symbol}
(x)_n := \begin{cases}
1, & \text{ if } n = 0,\\
x (x+1) \cdots (x+n-1), & \text{ if } n \geq 1.
\end{cases}
\end{equation}

\begin{equation}\label{kummer}
{}_3F_2\left(
\begin{matrix}
a, b, c\\
d, e
\end{matrix}\; ; 1\right) = \frac{\Gamma(e)\Gamma(d+e-a-b-c)}{\Gamma(e-a)\Gamma(d+e-b-c)}{}_3F_2\left(
\begin{matrix}
a, d-b, d-c\\
d, d+e-b-c
\end{matrix}\; ; 1\right).
\end{equation}

\begin{lemma}\label{lemma-P=Q}
Suppose $a \geq m+2N+2$. Then, the polynomials $P$ and $Q$ defined in \eqref{def-P(lambda)} and \eqref{def-Q(lambda)} above coincide.
\end{lemma}

\begin{proof}
In order to show $P = Q$, it suffices to prove the following three facts:
\begin{enumerate}[label=\normalfont{(\arabic*)}]
\item $\deg P = \deg Q = 2N+1$.
\item The top terms of $P$ and $Q$ coincide.
\item $P(\lambda) = 0$ for any $\lambda \in \Lambda(N,a) = \{N+1-a-q: q=0, 1, \ldots, 2N\}$. 
\end{enumerate}
The first one is straightforward, since $\alpha$ and $\beta$ have degree $N$ by \eqref{def-alpha} and \eqref{def-beta}, and since $Q$ is the product of $2N+1$ monomials.

The second one can be shown easily. In fact, by \eqref{definition-constants-Cjr+-}, the top term of $P$ is given by
\begin{equation*}
\left(\left[\frac{a+m}{2}\right]-\left[\frac{a-m}{2}\right]\right)\frac{\Gamma(N+m+1)\Gamma(N-m+1)}{\Gamma(1+m)\Gamma(1-m)},
\end{equation*}
which coincides with the top term of $Q$ since $\left[\frac{a+m}{2}\right]-\left[\frac{a-m}{2}\right] = m$.

Let us show (3), which is the hardest part. Set $\lambda = N+1-a-q$ with $q = 0, 1, \ldots, 2N$. Then, $\alpha, \beta$ amount to
\begin{align*}
\alpha(\lambda, N, a, m) &= \frac{\Gamma(2N-q+1)}{\Gamma(N+1)}\frac{\Gamma(\left[\frac{a+m}{2}\right])}{\Gamma(\left[\frac{a+m}{2}\right] -N + q)}\widetilde{\alpha}(a, m),\\
\beta(\lambda, N, a, m) &= \frac{\Gamma(2N-q+1)}{\Gamma(N+1)}\frac{\Gamma(\left[\frac{a+m+2}{2}\right])}{\Gamma(\left[\frac{a+m+2}{2}\right] -N + q)}\widetilde{\alpha}(a+2, m),
\end{align*}
where
\begin{equation*}
\widetilde{\alpha}(a, m) := \sum_{r=0}^{\min(q,N)}\begin{pmatrix}
N\\
r
\end{pmatrix} 
\frac{\Gamma(2N-r+1)}{\Gamma(2N-q+1)}\frac{\Gamma(q+1)}{\Gamma(q-r+1)}\frac{\Gamma(N-m+1)}{\Gamma(N-m-r+1)}\frac{\Gamma(\left[\frac{a+m}{2}\right] -N + q)}{\Gamma(\left[\frac{a+m}{2}\right] -N + r)}.
\end{equation*}
Now, by dividing $P(\lambda)$ by
\begin{equation}\label{proof-constant-dividing}
\left(\frac{\Gamma(2N-q+1)}{\Gamma(N+1)}\right)^2 \frac{\Gamma(\left[\frac{a-m+2}{2}\right])}{\Gamma(\left[\frac{a-m}{2}\right] -N + q)}\frac{\Gamma(\left[\frac{a+m+2}{2}\right])}{\Gamma(\left[\frac{a+m}{2}\right] -N + q)} \neq 0,
\end{equation}
we obtain that $P(\lambda) = 0$ if and only if 
\begin{equation}\label{proof-alpha-pm-equality}
\widetilde{\alpha}(a, m)\widetilde{\alpha}(a+2, -m) = \widetilde{\alpha}(a, -m)\widetilde{\alpha}(a+2, m).
\end{equation}
Let us show that we actually have
\begin{equation*}
\begin{aligned}
\widetilde{\alpha}(a, -m) &= \widetilde{\alpha}(a, m),\\
\widetilde{\alpha}(a+2, -m) &= \widetilde{\alpha}(a+2, m),
\end{aligned}
\end{equation*}
which in particular implies \eqref{proof-alpha-pm-equality}. Since $a$ is arbitrary, it suffices to show the first equality. By using the Pochhammer notation of the rising factorial (see \eqref{Pochhamer-symbol}), $\widetilde{\alpha}(a, \pm m)$ can be rewritten as
\begin{equation}\label{proof-alpha-hypergeometric}
\begin{aligned}
\widetilde{\alpha}(a, \pm m) &= (2N-q+1)_q \left(\left[\frac{a\pm m}{2}\right]-N\right)_q \sum_{r=0}^{\min(q,N)}\frac{(-q)_r(-N)_r( \pm m-N)_r}{(-2N)_r \left(\left[\frac{a\pm m}{2}\right]-N\right)_r r!}\\
& = (2N-q+1)_q \left(\left[\frac{a\pm m}{2}\right]-N\right)_q {}_3F_2\left(\begin{matrix}
-q, -N, \pm m -N\\
-2N, \left[\frac{a\pm m}{2}\right]-N
\end{matrix}\; ; 1\right),
\end{aligned}
\end{equation}
where ${}_3F_2$ denotes the hypergeometric series with $(p,q) = (3,2)$ (see \eqref{def-hypergeometric}). Thus, the equality $\widetilde{\alpha}(a, -m) = \widetilde{\alpha}(a, m)$ amounts to
\begin{equation*}
_3F_2\left(\begin{matrix}
-q, -N, -m -N\\
-2N, \left[\frac{a-m}{2}\right]-N
\end{matrix}\; ; 1\right) = \frac{\left(\left[\frac{a+ m}{2}\right]-N\right)_q}{\left(\left[\frac{a- m}{2}\right]-N\right)_q}{}_3F_2\left(\begin{matrix}
-q, -N, m -N\\
-2N, \left[\frac{a+ m}{2}\right]-N
\end{matrix}\; ; 1\right),
\end{equation*}
which holds now thanks to Kummer's identity \eqref{kummer}. Hence, \eqref{proof-alpha-pm-equality} holds; in other words, $P(\lambda) = 0$.
\end{proof}

\end{document}
